\textbf{Phần 1. Trắc nghiệm lựa chọn}
\setcounterpageref{ex}{0}

\begin{ex} %[2D11V56]
Mệnh giá tiền nào dưới đây KHÔNG có trong chương trình làm quen của học sinh lớp 2?
\choice
{\True $50 000\text{ đồng}$}
{$100\text{ đồng}$}
{$500\text{ đồng}$}
{$1 000\text{ đồng}$}
\loigiai{
Lớp 2 chỉ làm quen với: 100 đồng, 200 đồng, 500 đồng và 1 000 đồng.
\textbf{Đáp án của câu hỏi là A}
}
\end{ex}

\begin{ex} %[2D11V56]
Cần bao nhiêu tờ tiền $500\text{ đồng}$ để đổi được $1$ tờ tiền $1 000\text{ đồng}$?
\choice
{$5$ tờ}
{$3$ tờ}
{\True $2$ tờ}
{$10$ tờ}
\loigiai{
$500 \times 2 = 1 000$ đồng, nên cần 2 tờ.
\textbf{Đáp án của câu hỏi là C}
}
\end{ex}

\begin{ex} %[2D11V56]
Nam có $3$ tờ tiền $200\text{ đồng}$ và $1$ tờ tiền $100\text{ đồng}$. Nam có tất cả bao nhiêu đồng?
\choice
{$600\text{ đồng}$}
{\True $700\text{ đồng}$}
{$800\text{ đồng}$}
{$500\text{ đồng}$}
\loigiai{
$200 \times 3 + 100 = 700$ đồng.
\textbf{Đáp án của câu hỏi là B}
}
\end{ex}

\begin{ex} %[2D11V56]
Một cây bút chì giá $800\text{ đồng}$. Hoa đưa cho cô bán hàng tờ tiền $1 000\text{ đồng}$. Cô bán hàng trả lại bao nhiêu tiền thừa?
\choice
{$100\text{ đồng}$}
{$300\text{ đồng}$}
{$50\text{ đồng}$}
{\True $200\text{ đồng}$}
\loigiai{
$1 000 - 800 = 200$ đồng.
\textbf{Đáp án của câu hỏi là D}
}
\end{ex}

\begin{ex} %[2D11V56]
Có bao nhiêu tờ $200\text{ đồng}$ để đổi được đúng $1 000\text{ đồng}$?
\choice
{\True $5$ tờ}
{$4$ tờ}
{$6$ tờ}
{$2$ tờ}
\loigiai{
$1 000 : 200 = 5$ tờ.
\textbf{Đáp án của câu hỏi là A}
}
\end{ex}

\begin{ex} %[2D11V56]
An có $1$ tờ $500\text{ đồng}$ và $2$ tờ $200\text{ đồng}$. An mua một cục tẩy hết $600\text{ đồng}$. An còn lại bao nhiêu đồng?
\choice
{$200\text{ đồng}$}
{$400\text{ đồng}$}
{\True $300\text{ đồng}$}
{$100\text{ đồng}$}
\loigiai{
An có: $500 + 200 \times 2 = 900$ đồng. Còn lại: $900 - 600 = 300$ đồng.
\textbf{Đáp án của câu hỏi là C}
}
\end{ex}

\begin{ex} %[2D11V56]
Một chiếc kẹo giá $200\text{ đồng}$. Với $1 000\text{ đồng}$, em có thể mua được nhiều nhất bao nhiêu chiếc kẹo như thế?
\choice
{$4$ cái}
{\True $5$ cái}
{$6$ cái}
{$3$ cái}
\loigiai{
$1 000 : 200 = 5$ cái kẹo.
\textbf{Đáp án của câu hỏi là B}
}
\end{ex}

\begin{ex} %[2D11V56]
Tờ tiền $1 000\text{ đồng}$ có mệnh giá gấp bao nhiêu lần tờ tiền $100\text{ đồng}$?
\choice
{$5$ lần}
{$20$ lần}
{$2$ lần}
{\True $10$ lần}
\loigiai{
$1 000 : 100 = 10$ lần.
\textbf{Đáp án của câu hỏi là D}
}
\end{ex}

\begin{ex} %[2D11V56]
Bình có $1$ tờ $500\text{ đồng}$, $1$ tờ $200\text{ đồng}$ và $2$ tờ $100\text{ đồng}$. Tổng số tiền của Bình là
\choice
{\True $900\text{ đồng}$}
{$800\text{ đồng}$}
{$1 000\text{ đồng}$}
{$700\text{ đồng}$}
\loigiai{
$500 + 200 + 100 \times 2 = 900$ đồng.
\textbf{Đáp án của câu hỏi là A}
}
\end{ex}

\begin{ex} %[2D11V56]
Mua một quyển vở hết $700\text{ đồng}$, mua một cái thước kẻ hết $200\text{ đồng}$. Mua cả hai thứ hết bao nhiêu đồng?
\choice
{$800\text{ đồng}$}
{$1 000\text{ đồng}$}
{\True $900\text{ đồng}$}
{$600\text{ đồng}$}
\loigiai{
$700 + 200 = 900$ đồng.
\textbf{Đáp án của câu hỏi là C}
}
\end{ex}

\begin{ex} %[2D11V56]
Để có đúng $800\text{ đồng}$ từ các tờ tiền $200\text{ đồng}$, cần ít nhất bao nhiêu tờ $200\text{ đồng}$?
\choice
{$3$ tờ}
{\True $4$ tờ}
{$5$ tờ}
{$6$ tờ}
\loigiai{
$800 : 200 = 4$ tờ.
\textbf{Đáp án của câu hỏi là B}
}
\end{ex}

\begin{ex} %[2D11V56]
Một gói bánh giá $1 000\text{ đồng}$. Minh có $2$ tờ $200\text{ đồng}$ và $1$ tờ $500\text{ đồng}$. Minh còn thiếu bao nhiêu đồng nữa mới đủ mua?
\choice
{$200\text{ đồng}$}
{$50\text{ đồng}$}
{$300\text{ đồng}$}
{\True $100\text{ đồng}$}
\loigiai{
Minh có $400 + 500 = 900$ đồng. Thiếu: $1 000 - 900 = 100$ đồng.
\textbf{Đáp án của câu hỏi là D}
}
\end{ex}

\textbf{Phần 2. Trắc nghiệm đúng sai}
\setcounterpageref{ex}{0}

\begin{ex} %[2D11V56]
Kiểm tra các nhận định về các mệnh giá tiền Việt Nam:
\choiceTF
{\True 1 tờ 1 000 đồng đổi được 2 tờ 500 đồng.}
{\True 5 tờ 200 đồng có tổng giá trị bằng 1 000 đồng.}
{Tờ tiền 100 đồng có giá trị lớn hơn tờ tiền 500 đồng.}
{\True Có thể dùng 1 tờ 500 đồng và 2 tờ 200 đồng để được 900 đồng.}
\loigiai{
\begin{itemize}
    \item Ý a) Đúng: 500 + 500 = 1 000 đồng. Đúng.
    \item Ý b) Đúng: 200 x 5 = 1 000 đồng. Đúng.
    \item Ý c) Sai: 100 đồng bé hơn 500 đồng.
    \item Ý d) Đúng: 500 + 200 x 2 = 900 đồng. Đúng.
\end{itemize}
\begin{itemize}
    \item a) \textbf{Đúng}.
    \item b) \textbf{Đúng}.
    \item c) \textbf{Sai}.
    \item d) \textbf{Đúng}.
\end{itemize}
}
\end{ex}

\begin{ex} %[2D11V56]
Bác Ba đi chợ với 1 tờ 1 000 đồng, mua mớ rau giá 400 đồng:
\choiceTF
{\True Bác Ba còn lại 600 đồng tiền thừa.}
{\True Số tiền còn lại đổi được đúng 3 tờ 200 đồng.}
{Số tiền còn lại đổi được 2 tờ 500 đồng.}
{\True Nếu mua thêm 1 quả ớt giá 100 đồng thì bác còn lại 500 đồng.}
\loigiai{
\begin{itemize}
    \item Ý a) Đúng: 1 000 - 400 = 600 đồng. Đúng.
    \item Ý b) Đúng: 600 : 200 = 3 tờ. Đúng.
    \item Ý c) Sai: 2 tờ 500 đồng là 1 000 đồng > 600 đồng.
    \item Ý d) Đúng: 600 - 100 = 500 đồng. Đúng.
\end{itemize}
\begin{itemize}
    \item a) \textbf{Đúng}.
    \item b) \textbf{Đúng}.
    \item c) \textbf{Sai}.
    \item d) \textbf{Đúng}.
\end{itemize}
}
\end{ex}

\textbf{Phần 3. Trả lời ngắn}
\setcounterpageref{ex}{0}

\begin{ex} %[2D11V56]
Mai có 1 tờ 1 000 đồng. Mai mua 2 phong bì thư, mỗi phong bì giá 300 đồng. Mai còn lại bao nhiêu đồng tiền thừa?
\shortans{400}
\loigiai{
\textbf{Đáp số:} $400$\\
Mua hết: $300 \times 2 = 600$ đồng. Còn lại: $1 000 - 600 = 400$ đồng.
}
\end{ex}

\begin{ex} %[2D11V56]
Có 1 tờ 500 đồng và 2 tờ 200 đồng. Cần thêm bao nhiêu tờ 100 đồng nữa để đủ 1 000 đồng?
\shortans{1}
\loigiai{
\textbf{Đáp số:} $1$\\
Đã có 900 đồng. Thiếu: $(1 000 - 900) : 100 = 1$ tờ 100 đồng.
}
\end{ex}

\begin{ex} %[2D11V56]
Bác Ba có 10 tờ tiền 100 đồng. Bác muốn đổi sang loại tiền 500 đồng. Bác đổi được bao nhiêu tờ 500 đồng?
\shortans{2}
\loigiai{
\textbf{Đáp số:} $2$\\
$100 \times 10 = 1 000$ đồng. $1 000 : 500 = 2$ tờ.
}
\end{ex}

\begin{ex} %[2D11V56]
Một chiếc bánh rán giá 500 đồng. Hỏi 2 chiếc bánh rán giá bao nhiêu đồng?
\shortans{1 000}
\loigiai{
\textbf{Đáp số:} $1 000$\\
$500 \times 2 = 1 000$ đồng.
}
\end{ex}

\begin{ex} %[2D11V56]
Dũng có 4 tờ 200 đồng. Dũng có tất cả bao nhiêu đồng?
\shortans{800}
\loigiai{
\textbf{Đáp số:} $800$\\
$200 \times 4 = 800$ đồng.
}
\end{ex}

\begin{ex} %[2D11V56]
Huy có 1 000 đồng, Huy mua kem hết 700 đồng. Cô bán hàng thối lại tiền toàn bằng tờ 100 đồng. Hỏi cô đưa cho Huy mấy tờ 100 đồng?
\shortans{3}
\loigiai{
\textbf{Đáp số:} $3$\\
Thừa 300 đồng. $300 : 100 = 3$ tờ.
}
\end{ex}
